% dissipative.tex -- Methodology: the dissipative network and the tracker modes it uses
% (draft 4 Oct 2026). Code: dissipative_planner/dissipative_network.py (DissipativeParams,
% DissipativeNetwork.step / reference), dissipative_node.py (phase dispatch, hand-over),
% tracker/velocity_loop.py; design notes docs/design/velocity_loop.md (sections 3, 10, 11).

\subsection{Dissipative network}
\label{sec:dissipative}

The second reference generator is a port of the decentralised dissipative controller of Quan et
al.\ \cite{quan2025swarm}, adapted to rigid rods, a ring payload and the same trackers. It was built
to compare the two approaches of the literature review on the same plant (RQ2), and because a
member leaving or joining is its native mode of operation: the network needs no new solver and no
hand-designed redistribution when the team changes.

\paragraph{The network.} Every UAV is a node of a virtual spring--damper network, as in the
table-mechanics analogy of \cite{quan2025swarm}. Two further nodes are pinned. The payload node sits
at the desired payload position $\mathbf{p}_L^{\mathrm{des}}$, and the anchor node sits on the
vertical axis above it at the height $l \sin\phi_0$, with $\phi_0 = 45^\circ$ the design elevation
(the paper's fixed ``table top''). Each UAV node $i$ is joined by springs to the payload node (rest
length $l$), to the anchor (rest length $l\cos\phi_0$, the cone radius), to its two neighbours in
the ring (rest length the chord of an even polygon for the current number of members) and to its
azimuth slot, and is damped relative to all of them:
\begin{equation}
    m_n \ddot{\mathbf{x}}_i = \sum_{j \in \mathcal{N}_i} k_{ij}\big(\|\mathbf{x}_j - \mathbf{x}_i\| - L_{ij}\big)
    \frac{\mathbf{x}_j - \mathbf{x}_i}{\|\mathbf{x}_j - \mathbf{x}_i\|}
    \; - \; c \sum_{j \in \mathcal{N}_i} \big(\dot{\mathbf{x}}_i - \dot{\mathbf{x}}_j\big).
    \label{eq:network}
\end{equation}
The damping is the dissipation: any disturbance, including a member leaving, injects energy that it
removes. With every spring at its rest length the desired cone is a force-free equilibrium, and the
damping makes it the stable one. The network is integrated with semi-implicit Euler, ten substeps
per 10\,Hz tick. Its defaults are spring stiffnesses of 40, 40 and 20\,N/m to the payload, anchor
and neighbours, 18\,N/m to the slot, damping $c = 6$\,N\,s/m and a node mass $m_n = 0.5$\,kg, which
settle the network in about 1\,s.

\paragraph{From node to reference.} As in \cite{quan2025swarm}, the node position is the reference.
UAV $i$ is placed one rod length out from its own attachment point along the direction of its node,
\begin{equation}
    \mathbf{p}_i^{\mathrm{ref}} = \mathbf{a}_i + l\, \frac{\mathbf{x}_i - \mathbf{a}_i}{\|\mathbf{x}_i - \mathbf{a}_i\|},
    \qquad \mathbf{a}_i = \mathbf{p}_L^{\mathrm{des}} + R_\psi \boldsymbol{\rho}_i,
\end{equation}
where $R_\psi$ is the payload's measured yaw only. Every spring hangs off the node's own attachment
point on the rim, not the payload's centre. The cone stays vertical; tilting it onto the effective
gravity of the commanded acceleration, as the planner's references do, is one of the variants of
Section~\ref{sec:sim_architecture}.

\redo{the network's cone drawn from the payload's centre against from each UAV's own attachment point on the rim (offline network test): whether the rods can reach the geometry.}
% --- original text, flown before the rig twin; kept to refill this place:
% where $R_\psi$ is the payload's measured yaw only. Every spring hangs off the node's own attachment
% point on the rim, not the payload's centre. A cone drawn from the centre of a 0.5\,m ring asks for a
% geometry the rods cannot reach, and the payload fell in the offline tests. The cone stays vertical:
% tilting it onto the effective gravity of the commanded acceleration, as the planner's references do,
% was tried and gave no benefit (Section~\ref{sec:sim_architecture}).

\paragraph{Cable feedforward.} The trackers need the rod's pull, which the paper's velocity-commanded
robots do not. With $m$ UAVs attached, each gets an equal share
$\lambda_i = m_L g / (m \sin\phi_i)$ at its node's elevation $\phi_i$. For an uneven layout, where
equal shares would leave a net moment and tilt the payload, the tensions instead come from a
six-degree-of-freedom wrench balance at the actual attachment geometry.

\paragraph{Open loop in payload pose.} The network is built on the desired payload position and the
measured yaw only. It does not use the payload's measured position or roll and pitch. Building the
geometry or the tensions on the measured attitude, or adding a horizontal leash on the measured
position, closes a loop on the payload's pose; these variants are in
Section~\ref{sec:sim_architecture}. The rods close the loop physically.
% --- original text, flown before the rig twin; kept to refill this place:
% measured yaw only. It does not use the payload's measured position or roll and pitch. Building the
% geometry or the tensions on the measured attitude, or adding a horizontal leash on the measured
% position, closes a loop that amplifies tilt: each variant tilted the payload further or aborted the
% flight (Section~\ref{sec:sim_architecture}). The rods close the loop physically.

\paragraph{Common-mode trim.} The tracker has no integrator, so a steady thrust error shows as a
steady sag. One integrator on the measured payload height error adds the same correction to every
attached UAV's acceleration feedforward, bounded to 2\,m/s$^2$:
\begin{equation}
    \mathbf{a}_{\mathrm{trim}} = k_{\mathrm{trim}} \int \big(z_L^{\mathrm{des}} - z_L\big)\, dt\; \mathbf{e}_3 .
\end{equation}
Per-UAV integrators are not used: each UAV integrates its own
error, and the differences become a moment (Section~\ref{sec:sim_architecture}). A common-mode trim cannot produce a differential moment.
% --- original text, flown before the rig twin; kept to refill this place:
% Per-UAV integrators were tried first and tilted the payload to 60$^\circ$: each UAV integrates its own
% error, and the differences become a moment. A common-mode trim cannot produce a differential moment.
A variant scales the trim by each UAV's tension share, so that on an uneven layout the most heavily
loaded UAV, which sags most, gets the most correction.

\paragraph{Velocity tracking.} The robots of \cite{quan2025swarm} take a velocity command, and their
autopilot closes velocity, attitude and rate. Betaflight accepts only body rates and throttle, so the
velocity layer is supplied by an alternative tracker mode at 50\,Hz:
\begin{equation}
    \mathbf{v}_{\mathrm{sp}} = \mathbf{v}^{\mathrm{ref}} + k_p\big(\mathbf{p}^{\mathrm{ref}} - \mathbf{p}\big),
    \qquad
    \mathbf{a}_{\mathrm{sp}} = \mathbf{a}^{\mathrm{ff}} + k_v\big(\mathbf{v}_{\mathrm{sp}} - \mathbf{v}\big),
\end{equation}
from which the throttle and the attitude set-point follow, and the rate command follows from the
attitude error through the inverse of the flight controller's rate curve. This loop cannot fly the
lift, so the mode used with the network flies the MPC tracker through the lift and switches to the
velocity loop when the network takes over.

\paragraph{Switching between planner and network.} The planner always flies the creep, the lift and
the carry; the network cannot break the payload off the floor. The network takes over on a detach
(in the network detach mode) or by command. At the switch its desired payload position is set to the
measured one, so there is no jump in the references, and the cable feedforward blends from the
planner's values to the network's over a ramp. After a detach the network holds the payload while it
redistributes the load, and then hands back to a planner solver for the smaller team, primed twice
from the measured state before it flies.
